%%%%%%%%%%%%%%%%%%%%%%%%%%%%%%%%%%%%%%%%%%%%%%%%%%
% Metadata
%%%%%%%%%%%%%%%%%%%%%%%%%%%%%%%%%%%%%%%%%%%%%%%%%%
% id: jmo2026y1
% title: 2026年 第36回 日本数学オリンピック 予選問題 第1問
% tags: [整数]
% difficulty: 2
% source: https://www.imojp.org/archive/mo2026/jmo/problems/jmo36yqa.pdf

%%%%%%%%%%%%%%%%%%%%%%%%%%%%%%%%%%%%%%%%%%%%%%%%%%
% Preamble
%%%%%%%%%%%%%%%%%%%%%%%%%%%%%%%%%%%%%%%%%%%%%%%%%%
\documentclass[fleqn]{ltjsarticle}

% パッケージ
\usepackage{luatexja}
\usepackage{amsmath,amssymb,amsfonts}
\usepackage{common,ascolorbox}
\usepackage[ceo]{emath}

% 行間や文字間隔の調整
\thickmuskip=1.0\thickmuskip
\medmuskip=0.8\medmuskip
\thinmuskip=0.8\thinmuskip
\arraycolsep=0.3\arraycolsep
\AtBeginDocument{
  \abovedisplayskip=0.85\abovedisplayskip
  \belowdisplayskip=0.55\belowdisplayskip
  \abovedisplayshortskip=0.5\abovedisplayshortskip
  \belowdisplayshortskip=0.5\belowdisplayshortskip
}

% ヘッダー
\usepackage{fancyhdr}
\pagestyle{fancy}
\renewcommand{\headrulewidth}{0pt}
\lhead{\textbf{2026年 第36回 日本数学オリンピック 予選問題 第1問}}
\rhead{\repeatstars{2}}
\cfoot{}

%%%%%%%%%%%%%%%%%%%%%%%%%%%%%%%%%%%%%%%%%%%%%%%%%%
% Document
%%%%%%%%%%%%%%%%%%%%%%%%%%%%%%%%%%%%%%%%%%%%%%%%%%
\begin{document}

\begin{problembox}{問題}
  $a^{20}+b^{2}+c^{6}=2026$ をみたす正の整数の組 $(a, b, c)$ をすべて求めよ．
\end{problembox}

\noindent
\kaib\quad
$a\geqq 2$のとき
\begin{align*}
    a^{20}+b^{2}+c^{6}>a^{20}>2026
\end{align*}
であるから，$a=1$ である必要がある．$a=1$ を代入して
\begin{align*}
    b^{2}+c^{6}=2025 \Y b^{2}=2025-c^{6}=(45+c^{3})(45-c^{3})
\end{align*}
をみたす正の整数の組 $(b, c)$ を求める． \\
$b^{2}>0, 45+c^{3}>0$ から $45-c^{3}>0$ である必要があるので，$c$ は $1, 2, 3$ のいずれかである． \\
それぞれ代入して計算すると
\begin{alignat*}{2}
    & b^{2}=(45+1^{3})(45-1^{3})=46\times 44 & & (c=1) \\
    & b^{2}=(45+2^{3})(45-2^{3})=53\times 37 & & (c=2) \\
    & b^{2}=(45+3^{3})(45-3^{3})=72\times 18=36^{2} & \quad & (c=3)
\end{alignat*}
より，$c=3$ のとき $b=36$ を得る． \\
以上から
\begin{align*}
    \boldsymbol{(a, b, c)=(1, 36, 3)}\kotae
\end{align*}

% \noindent
% \sankoua\quad
% 特に書くことなし

\end{document}